\chapter{Foundations of Architecture}

\section{The Concept of Tokens}
Computers cannot understand raw text. They only understand numbers. \textbf{Tokenization} is the process of converting text into a sequence of integers.

\begin{tcolorbox}[colback=blue!5!white,colframe=blue!75!black,title=Analogy: The Language of Numbers]
Imagine trying to teach a computer to read a book. You can't just show it the pages. Instead, you assign a unique ID number to every word or sub-word (like "ing" or "pre"). The computer then reads a stream of IDs like `[101, 45, 2099]`. This sequence of numbers is the input to our model.
\end{tcolorbox}

In `pure_transformer`, we use a tokenizer (like GPT-2's Byte-Pair Encoding) to handle this conversion.

\section{Embeddings: Giving Meaning to Numbers}
A simple integer ID like `2099` implies order (2099 > 101), which is meaningless for words. To fix this, we use \textbf{Embeddings}.
An embedding is a dense vector (a list of floating-point numbers) that represents the \textit{meaning} of a token.
\begin{itemize}
    \item King: `[0.9, 0.1, ...]`
    \item Queen: `[0.9, 0.2, ...]`
\end{itemize}
Ideally, words with similar meanings will have similar embedding vectors.

\section{The High-Level Data Flow}
In a decoder-only transformer, the flow of data is straightforward:
\begin{enumerate}
    \item \textbf{Input Text:} "The cat sat on the"
    \item \textbf{Tokenization:} Converted to integers (IDs): \texttt{[464, 3797, 3373, 319, 262]}
    \item \textbf{Embedding:} Integers are mapped to dense vectors (e.g., size 1024).
    \item \textbf{Layers:} Vectors pass through $N$ transformer blocks.
    \item \textbf{Output Head:} The final vectors are projected back to vocabulary size.
    \item \textbf{Softmax:} Probabilities for the next token (e.g., "mat").
\end{enumerate}

\section{Embeddings}
The embedding layer serves as the look-up table. In `pure\_transformer`, we use PyTorch's `nn.Embedding`.

\begin{lstlisting}[language=Python, caption=Embedding Layer in TransformerLM]
# model/transformer.py

class TransformerLM(nn.Module):
    def __init__(self, config: TransformerConfig):
        super().__init__()
        # ...
        # Token embeddings
        self.embed_tokens = nn.Embedding(config.vocab_size, config.hidden_size)
\end{lstlisting}

This maps a discrete token ID (from 0 to 50,303) to a continuous vector of size `hidden_size` (e.g., 1024).

\section{Positional Encodings: RoPE}
Transformers process all tokens in parallel, so they have no inherent notion of "order". We must inject positional information.

Early models added learnable vectors to the embeddings. Modern models use \textbf{Rotary Positional Embeddings (RoPE)}. RoPE rotates the query and key vectors in the attention mechanism, which mathematically encodes the relative distance between tokens.

\subsection{Implementation}
RoPE is efficient because tokens farther apart have a smaller dot product (due to the rotation), naturally decaying the attention score.

\begin{lstlisting}[language=Python, caption=Precomputing RoPE Cache]
# model/transformer.py

def precompute_rope_cache(
    seq_len: int,
    head_dim: int,
    theta: float = 10000.0,
    # ...
) -> Tuple[Tensor, Tensor]:
    """Precompute RoPE cos/sin cache."""
    # Frequencies decay exponentially
    inv_freq = 1.0 / (theta ** (torch.arange(0, head_dim, 2) / head_dim))
    t = torch.arange(seq_len)
    freqs = torch.outer(t, inv_freq)
    
    # Cosine and Sine components for rotation
    cos = freqs.cos()
    sin = freqs.sin()
    return cos, sin
\end{lstlisting}

These precomputed cosine and sine values are then applied to the Query (Q) and Key (K) vectors during the attention step.
